\documentclass{article}
\usepackage[utf8]{inputenc}
\usepackage{amsmath}
\usepackage{amssymb}
\usepackage{geometry}
\geometry{margin=1in}

\title{Qubit Electronics-Formula Sheet}
\author{Martina Paola Calvi}
\date{}

\begin{document}

\maketitle
\section{Qubit Introduction}
\subsection*{Fundamental Qubit Representation and State Vectors}
\begin{itemize}
    \item \textbf{General Superposition:} $|\psi\rangle = \alpha|0\rangle + \beta|1\rangle = \begin{bmatrix} \alpha \\ \beta \end{bmatrix}$ where $\alpha$ and $\beta$ are complex probability amplitudes.
    \item \textbf{Normalization Condition:} $|\alpha|^2 + |\beta|^2 = 1=\langle\psi|\psi\rangle$
    \item \textbf{Hermitian Conjugate:} $|\psi\rangle^\dagger = \langle\psi| = [\alpha^*, \beta^*]$

    \item \textbf{Multi-Qubit System Scaling:} A system of $N$ qubits scales to $2^N$ components. For $N=3$: $|\psi\rangle = c_1|000\rangle + c_2|001\rangle + \cdots + c_8|111\rangle$
\end{itemize}

\subsection*{The Bloch Sphere Geometry}
    \begin{itemize}
    \item \textbf{General State Formula:} $|\psi\rangle = \cos\left(\frac{\theta}{2}\right)|0\rangle + e^{i\phi}\sin\left(\frac{\theta}{2}\right)|1\rangle$
    \item \textbf{Bloch Vector Coordinates:} $|0\rangle$ (North Pole): $(0, 0, 1)$; $|1\rangle$ (South Pole): $(0, 0, -1)$; $|+\rangle$: $(1, 0, 0)$; $|-\rangle$: $(-1, 0, 0)$
\end{itemize}

\subsection*{Quantum Measurement and the Born Rule}
\begin{itemize}
    \item \textbf{State Collapse:} Measurement forces $|\psi\rangle$ to collapse into either $|0\rangle$ or $|1\rangle$.
    \item \textbf{Born Rule:} $P(|a\rangle) = |\langle \psi|a \rangle|^2$
\end{itemize}

\subsection*{Basis States in Different Representations}
\begin{itemize}
    \item \textbf{Computational Basis:} $|0\rangle = \begin{bmatrix} 1 \\ 0 \end{bmatrix}$, $|1\rangle = \begin{bmatrix} 0 \\ 1 \end{bmatrix}$
    \item \textbf{Hadamard Basis:} $|+\rangle = \frac{1}{\sqrt{2}}(|0\rangle + |1\rangle)$, $|-\rangle = \frac{1}{\sqrt{2}}(|0\rangle - |1\rangle)$
    \item \textbf{Circular Basis:} $|+i\rangle = \frac{1}{\sqrt{2}}(|0\rangle + i|1\rangle)$, $|-i\rangle = \frac{1}{\sqrt{2}}(|0\rangle - i|1\rangle)$
\end{itemize}

\subsection*{Single-Qubit Unitary Operators}
\begin{itemize}
    \item \textbf{X (NOT):} $\begin{bmatrix} 0 & 1 \\ 1 & 0 \end{bmatrix}$
    \item \textbf{Y:} $\begin{bmatrix} 0 & -i \\ i & 0 \end{bmatrix}$
    \item \textbf{Z:} $\begin{bmatrix} 1 & 0 \\ 0 & -1 \end{bmatrix}$
    \item \textbf{Hadamard (H):} $\frac{1}{\sqrt{2}}\begin{bmatrix} 1 & 1 \\ 1 & -1 \end{bmatrix}$
    \item \textbf{Phase (S):} $\begin{bmatrix} 1 & 0 \\ 0 & i \end{bmatrix}$
    \item \textbf{T:} $\begin{bmatrix} 1 & 0 \\ 0 & e^{i\pi/4} \end{bmatrix}$
    \item \textbf{Rotation Gates:} $R_x(\theta) = e^{-i\theta X/2}, \quad R_y(\theta) = e^{-i\theta Y/2}, \quad R_z(\theta) = e^{-i\theta Z/2}$
\end{itemize}

\subsection*{Two-Qubit Interactions: CNOT Gate}
\begin{itemize}
    \item \textbf{Mapping Rule:} $|A\rangle|B\rangle \to |A\rangle|A \oplus B\rangle$
    \item \textbf{Matrix Representation:} $\text{CNOT} = \begin{bmatrix} 1 & 0 & 0 & 0 \\ 0 & 1 & 0 & 0 \\ 0 & 0 & 0 & 1 \\ 0 & 0 & 1 & 0 \end{bmatrix}$
\end{itemize}

\subsection*{Density Matrix Formalism}
\begin{itemize}
    \item \textbf{Definition:} $\rho = |\psi\rangle\langle\psi|$
    \item \textbf{Mixed State Ensemble:} $\widehat{\rho}  \equiv  \sum_i p_i|\Psi_i\rangle\langle\Psi_i|$ where $p_i$ are probabilities and $\sum_i p_i = 1$
    \item \textbf{Trace Property:} $\text{Tr}(\rho^2) = 1 $ if pure state;$ < 1$ if mixed state
\end{itemize}

\subsection*{Qubit Performance Metrics}
\begin{itemize}
    \item \textbf{Coherence Time:} $\frac{1}{T_2} = \frac{1}{T_\phi} + \frac{1}{2T_1}$
    \item \textbf{Gate Capacity:} $N_{\text{gates}} \sim \frac{T_2}{t_{\text{gate}}}$
\end{itemize}


\newpage
\section{Superconducting Qubits}

\subsection*{Fundamental Superconductivity and Cooper Pairs}
\begin{itemize}
    \item \textbf{Collective Wave Function:} $\psi = \psi_0 \exp\left[-\mathrm{j}\frac{1}{\hbar}(Et - \mathbf{p} \cdot \mathbf{x})\right]$ where $\mathrm{j} = \sqrt{-1}$
    \item \textbf{Cooper Pair Density:} $\rho = \psi \cdot \psi^* = |\psi_0|^2$
    \item \textbf{Meissner Effect:} $B = 0$
\end{itemize}

\subsection*{Josephson Junction Dynamics}
\begin{itemize}
    \item \textbf{Phase Difference:} $\phi = \varphi_2 - \varphi_1$
    \item \textbf{Current-Phase Relation:} $I = I_C \sin \phi$
    \item \textbf{Voltage-Phase Relation:} $V = \frac{\Phi_0}{2\pi}\dot{\phi}$ where $\Phi_0 = \frac{h}{2e} = \frac{2\pi\hbar}{2e}$
    \item \textbf{Josephson Inductance:} $L = \frac{\Phi_0}{2\pi I_C \cos \phi}$ and $L_j = \frac{\Phi_0}{2\pi I_C}$
    \item \textbf{Josephson Energy:} $E_j(1 - \cos \phi)$
\end{itemize}

\subsection*{LC Circuit Quantization}
\begin{itemize}
    \item \textbf{Hamiltonian:} $H = \frac{Q^2}{2C} + \frac{\Phi^2}{2L}$
    \item \textbf{Resonant Frequency:} $\omega_{01} = \frac{1}{\sqrt{LC}}$
    \item \textbf{Canonical Commutation:} $[\hat{\Phi}, \hat{Q}] = \mathrm{i}\hbar$ and $[\hat{a}, \hat{a}^\dagger] = 1$
\end{itemize}

\subsection*{Transmon Qubit Parameters}
\begin{itemize}
    \item \textbf{Non-linear Inductance:} $L(\Phi) = \frac{\Phi_0}{2\pi I_C \cos(2\pi\Phi/\Phi_0)}$
    \item \textbf{Transition Energy:} $E_{01} = \hbar\omega_{01} \approx 10 \text{ GHz}$
    \item \textbf{Anharmonicity:} $E_{01} \neq E_{12}$ allows selective qubit control
    \item \textbf{Operating Temperature:} $\approx 10 \text{ mK}$, requiring $k_B T \ll \hbar\omega_{01}$
    \item \textbf{Gate Time:} $T_\pi \approx 4$--$20 \text{ ns}$
    \item \textbf{Coherence Times:} $T_1, T_2 \sim 100$--$1000 \text{ } \mu\text{s}$
\end{itemize}

\newpage
\section{Circuit QED}
\subsection*{Fundamental Light-Matter Interaction}
\begin{itemize}
    \item \textbf{Dipole Coupling:} $H = -\vec{d} \cdot \vec{E}(t)$ where $\vec{d}$ is the electric dipole moment
    \item \textbf{Rabi Frequency:} $\Omega_R = \frac{d E_0}{\hbar}$ for classical field amplitude $E_0$
\end{itemize}

\subsection*{Resonator Physics}
\begin{itemize}
    \item \textbf{Fabry-Pérot Resonance:} $2L = m\lambda_m$ where $m$ is the mode number
    \item \textbf{Mode Separation:} $\Delta f = \frac{c}{2L}$
    \item \textbf{LC Resonator Frequency:} $\omega_r = \frac{1}{\sqrt{L_r C_r}}$
    \item \textbf{Quality Factor:} $Q = \frac{\omega_r}{\kappa}$ where $\kappa$ is the decay rate
\end{itemize}

\subsection*{Jaynes-Cummings Model}
\begin{itemize}
    \item \textbf{Full Hamiltonian:} $H = \hbar\omega_r\hat{a}^\dagger\hat{a} + \frac{\hbar\omega_{01}}{2}\hat{\sigma}_z + \hbar g(\hat{a}\hat{\sigma}_+ + \hat{a}^\dagger\hat{\sigma}_-)$
    \item \textbf{Rotating Wave Approximation (RWA):} Neglects rapidly oscillating terms at frequency $\omega_r + \omega_{01}$
    \item \textbf{Jaynes-Cummings (RWA):} $H_{JC} = \hbar\omega_r\hat{a}^\dagger\hat{a} + \frac{\hbar\omega_{01}}{2}\hat{\sigma}_z + \hbar g(\hat{a}\hat{\sigma}_+ + \hat{a}^\dagger\hat{\sigma}_-)$
\end{itemize}

\subsection*{Coupling Strength}
\begin{itemize}
    \item \textbf{Coupling Formula:} $g = \omega_r \frac{C_g}{C} \left(\frac{E_j}{2E_c}\right)^{1/4} \sqrt{\frac{\pi Z_r}{R_k}}$ where $C_g$ is gate capacitance, $Z_r$ is resonator impedance, and $R_k = \frac{h}{e^2} \approx 25.8 \text{ k}\Omega$
    \item \textbf{Typical Range:} $g/2\pi \approx 100$ MHz to $1$ GHz (strong coupling regime)
\end{itemize}

\subsection*{Dispersive Regime}
\begin{itemize}
    \item \textbf{Condition:} $|\Delta| = |\omega_{01} - \omega_r| \gg g$
    \item \textbf{Dispersive Shift (AC Stark Shift):} $\chi = \frac{g^2}{\Delta}$
    \item \textbf{Effective Hamiltonian:} $H \approx (\omega_r + \chi \hat{\sigma}_z) \hat{a}^\dagger \hat{a} + \frac{\omega_{01}}{2} \hat{\sigma}_z$
    \item \textbf{State-Dependent Resonator Frequency:} $\omega_r^{(0)} = \omega_r - \chi$ (ground state $|0\rangle$); $\omega_r^{(1)} = \omega_r + \chi$ (excited state $|1\rangle$)
    \item \textbf{Readout via Dispersive Shift:} Phase accumulation $\propto \chi \hat{\sigma}_z t$ enables qubit-state discrimination without excitation
\end{itemize}

\subsection*{Circuit QED Implementation}
\begin{itemize}
    \item \textbf{Cavity Analogue:} Superconducting coplanar waveguide (CPW) replaces optical Fabry-Pérot mirrors
    \item \textbf{Artificial Atom:} Transmon qubit (Duffing oscillator) coupled to resonator via capacitive coupling $C_g$
    \item \textbf{Substrate:} Aluminum or niobium on sapphire for low loss
    \item \textbf{Field Confinement:} Tightly confined to $\sim 10^{-3}\lambda$, enabling large vacuum fluctuations $E_0 \approx 0.2 \text{ V/m}$
    \item \textbf{Coupling Enhancement:} Circuit QED achieves $\approx 2 \times 10^4 \times$ stronger coupling than cavity QED
\end{itemize}

\subsection*{System Hamiltonian Hierarchy}
\begin{itemize}
    \item \textbf{Full Circuit Ham.:} $H = \frac{Q_r^2}{2C_r} + \frac{\Phi_r^2}{2L_r} + 4E_c \hat{n}^2 - E_j \cos \hat{\phi} - g(\hat{b}^\dagger - \hat{b})(\hat{a}^\dagger - \hat{a})$
    \item \textbf{Two-Level Approx.:} $H = \omega_r \hat{a}^\dagger \hat{a} + \frac{\omega_{01}}{2} \hat{\sigma}_z - g(\hat{\sigma}_+ + \hat{\sigma}_-)(\hat{a}^\dagger + \hat{a})$
    \item \textbf{After RWA:} $H = \omega_r \hat{a}^\dagger \hat{a} + \frac{\omega_{01}}{2} \hat{\sigma}_z + g(\hat{a}^\dagger \hat{\sigma}_- + \hat{a} \hat{\sigma}_+)$
\end{itemize}

\subsection*{Transmon-Specific Considerations}
\begin{itemize}
    \item \textbf{Anharmonicity:} $\delta = (E_{12} - E_{01})/\hbar$ from Duffing nonlinearity; enables selective single-qubit addressing
    \item \textbf{Flux Tunability:} SQUID integration allows $\omega_{01}(\Phi_{\text{ext}})$ tuning via external magnetic flux
    \item \textbf{Charging Energy:} $E_c = \frac{e^2}{2C}$ determines qubit-cavity detuning range
    \item \textbf{Josephson Energy:} $E_j$ sets nonlinearity strength; transmon regime requires $E_j \gg E_c$
\end{itemize}

\newpage
\section{Qubit Control}
\subsection*{Readout and the Dispersive Regime}
\begin{itemize}
    \item \textbf{Full Jaynes-Cummings Hamiltonian:} $\hat{H} = \omega_r \hat{a}^\dagger \hat{a} + \frac{\omega_q}{2} \hat{\sigma}_z - g(\hat{a}^\dagger \hat{\sigma}_- + \hat{a} \hat{\sigma}_+)$
    \item \textbf{Dispersive Approximation:} $\hat{H} \approx (\omega_r + \chi \hat{\sigma}_z) \hat{a}^\dagger \hat{a} + \frac{\omega_q}{2} \hat{\sigma}_z$
    \item \textbf{Dispersive Regime Condition:} $g \ll |\omega_q - \omega_r|$, where $\Delta = \omega_q - \omega_r \gg g, \kappa$
    \item \textbf{Purcell Decay Rate:} $\gamma_p = \left(\frac{g}{\Delta}\right)^2 \kappa$
    \item \textbf{Cavity Decay Rate:} $\kappa = \frac{\omega_r}{Q}$
    \item \textbf{Optimal Readout Balance:} $\chi/\kappa \approx 1$
    \item \textbf{Purcell Filter:} Frequency-selective element with $\kappa_{\text{eff}}(\omega_q) \ll \kappa(\omega_r)$
\end{itemize}

\subsection*{Transmon Hamiltonian and Microwave Drive}
\begin{itemize}
    \item \textbf{Circuit Hamiltonian:} $\hat{H} = \frac{\hat{Q}^2}{2C_\Sigma} + \frac{\hat{\Phi}^2}{2L} + \frac{C_d}{C_\Sigma} V_d(t) \hat{Q}$
    \item \textbf{Charge Operator:} $\hat{Q} = -i Q_{\text{zpf}}(\hat{a} - \hat{a}^\dagger)$ where $Q_{\text{zpf}} = \sqrt{\frac{\hbar}{2Z_0}}$
    \item \textbf{Characteristic Impedance:} $Z_0 = \sqrt{\frac{L}{C_\Sigma}}$
    \item \textbf{Two-Level Truncation:} $\hat{a} \to \hat{\sigma}_-$, $\hat{a}^\dagger \to \hat{\sigma}_+$
    \item \textbf{Truncated Qubit Hamiltonian:} $\hat{H} = -\frac{\omega_q}{2} \hat{\sigma}_z + \Omega V_d(t) \hat{\sigma}_y$ where $\Omega = \frac{C_d}{C_\Sigma} Q_{\text{zpf}}$
\end{itemize}

\subsection*{Rotating Frame and Rotating Wave Approximation}
\begin{itemize}
    \item \textbf{Unitary Transformation:} $U(t) = e^{i\omega_q t \hat{\sigma}_z / 2}$
    \item \textbf{Monochromatic Drive:} $V_d(t) = V_0 s(t) (I \sin(\omega_d t) - Q \cos(\omega_d t))$
    \item \textbf{Effective Drive Hamiltonian (RWA):} $\hat{H}_d' = -\frac{\Omega}{2} V_0 s(t) (I\hat{\sigma}_x + Q\hat{\sigma}_y)$
    \item \textbf{RWA Validity:} Neglects rapidly oscillating terms at $\omega_d + \omega_q$; valid when $\omega_d \approx \omega_q$
\end{itemize}

\subsection*{Virtual Z Gates and Phase Control}
\begin{itemize}
    \item \textbf{Software Phase Update:} $I' = I \cos \phi - Q \sin \phi$; $Q' = I \sin \phi + Q \cos \phi$
    \item \textbf{Example (Virtual Z, $\phi = \pi/2$):} Initial pulse $(I=1, Q=0) \to (I'=0, Q'=1)$ converts X-rotation to Y-rotation
\end{itemize}


\section{Qubit Characterization}

\subsection*{Bloch Sphere and State Representation}
\begin{itemize}
    \item \textbf{Qubit State:} $|\psi\rangle = \cos\left(\frac{\theta}{2}\right)|0\rangle + e^{i\phi}\sin\left(\frac{\theta}{2}\right)|1\rangle$
    \item \textbf{Bloch Vector:} $\vec{r} = \begin{pmatrix} r_x \\ r_y \\ r_z \end{pmatrix} = \begin{pmatrix} \sin\theta\cos\phi \\ \sin\theta\sin\phi \\ \cos\theta \end{pmatrix}$
\end{itemize}

\begin{center}
\begin{tabular}{lll}
\textbf{Region} & \textbf{Axis Designation} & \textbf{Physical Counterpart} \\
\hline
North Pole & $+z$ & Longitudinal Axis \\
South Pole & $-z$ & Longitudinal Axis \\
Equator & $x$-$y$ Plane & Transverse Plane \\
Precession & $\phi(t)$ & Rotation at $\omega_{01}$ (Lab Frame)
\end{tabular}
\end{center}

\subsection*{Reference Frames and Resonance}
\begin{itemize}
    \item \textbf{Laboratory Frame:} The physical frame where the Bloch vector precesses around the $z$-axis at the qubit transition frequency $\omega_{01}$ when a transverse component exists.
    \item \textbf{Rotating Frame:} A reference frame rotating at frequency $\omega_r$. When $\omega_r = \omega_{01}$, the Bloch vector appears stationary.
    \item \textbf{Resonance Condition:} On resonance in the rotating frame, the ideal interaction Hamiltonian $H'$ is zero. Any observed movement in this frame is due to drive fields or noise.
\end{itemize}

\subsection*{Decoherence and Relaxation Rates}
\begin{itemize}
    \item \textbf{Energy Relaxation ($T_1$):} $\Gamma_1 \equiv 1/T_1$
    \item \textbf{Rate Synthesis:} $\Gamma_1 = \Gamma_{1\downarrow} + \Gamma_{1\uparrow}$
    \item \textbf{Boltzmann Excitation:} $\Gamma_{1\uparrow} = \Gamma_{1\downarrow}e^{-\hbar\omega/k_B T}$
    \item \textbf{Pure Dephasing:} $\Gamma_\phi$ arises from longitudinal noise coupling via the $z$-axis and causes frequency fluctuations $\delta\omega(t)$.
    \item \textbf{Transverse Relaxation ($T_2^*$):} $\Gamma_2 = \frac{\Gamma_1}{2} + \Gamma_\phi \quad \text{or} \quad \frac{1}{T_2^*} = \frac{1}{2T_1} + \frac{1}{T_\phi}$
    \item \textbf{Phase-Breaking Justification:} Energy relaxation is phase-breaking; once $|1\rangle \rightarrow |0\rangle$, phase information is destroyed.
    \item \textbf{Technology Note:} Transmon qubits are generally $T_1$-limited and typically exhibit relatively large $T_2$ times.
\end{itemize}

\subsection*{Characterization Measurement Protocols}
\begin{itemize}
    \item \textbf{Time of Flight (TOF):} Calibrates the propagation delay between readout pulse generation and signal arrival at the acquisition IP within the RFSoC.
    \item \textbf{Requirement:} Mandatory first step; applied to all subsequent measurements.
    \item \textbf{Resonator Spectroscopy:} Sweeps the readout tone to find the fundamental frequency; high power is used to decouple the qubit from the resonator.
    \item \textbf{Drive Spectroscopy:} Two-tone experiment used to calibrate the drive frequency to the qubit transition frequency $f_q$.
    \item \textbf{Calibration Example:} Identified $f_q = 3774.16$ MHz.
    \item \textbf{Rabi Oscillations:} Used to calibrate frequency, duration, and gain after identifying $f_q$.
    \item \textbf{Calibration Example (Digital Gain):} Calibrated $\pi$-pulse: $0.054$; calibrated $\pi/2$-pulse: $0.027$.
    \item \textbf{Resonator Punchout:} 2D sweep of readout gain versus frequency to map power response and identify the dressed readout frequency.
    \item \textbf{Physics:} At high power, the Josephson junction nonlinearity is suppressed; as power decreases, the resonance shifts to the dressed/dispersive readout frequency.
\end{itemize}

\subsection*{Coherence Measurements}
\begin{itemize}
    \item \textbf{$T_1$ Measurement:} $f(t) \propto e^{-t/T_1}$; signal decay of 63\% occurs at $t=T_1$.
    \item \textbf{Ramsey Protocol ($T_2^*$):} $X_{\pi/2}$ --- wait $\tau$ --- $X_{\pi/2}$ --- measure.
    \item \textbf{Mechanism:} Drive is detuned by $\delta\omega$ to induce controlled precession.
    \item \textbf{Spin Echo ($T_2$):} $X_{\pi/2}$ --- wait $\tau/2$ --- $Y_\pi$ --- wait $\tau/2$ --- $X_{\pi/2}$ --- measure.
    \item \textbf{Purpose:} The central $\pi$-pulse compensates for quasi-static noise.
    \item \textbf{Interpretation:} $T_2$ (echo) $>$ $T_2^*$ (Ramsey) indicates low-frequency noise.
\end{itemize}


\section{TLine}

\subsection*{Transmission Line Parameters}
\begin{itemize}
    \item \textbf{Resistance per unit length:} $R_0 \; [\Omega/\text{m}]$
    \item \textbf{Inductance per unit length:} $L_0 \; [\text{H}/\text{m}]$
    \item \textbf{Capacitance per unit length:} $C_0 \; [\text{F}/\text{m}]$
    \item \textbf{Conductance per unit length:} $G_0 \; [\text{S}/\text{m}]$
\end{itemize}

\subsection*{The Telegrapher Equations}
\begin{itemize}
    \item \textbf{Voltage equation:} 
    \[
    \frac{\partial V(z,t)}{\partial z} = -R_0\, i(z,t) - L_0 \frac{\partial i(z,t)}{\partial t}
    \]
    \item \textbf{Current equation:}
    \[
    \frac{\partial i(z,t)}{\partial z} = -G_0\, V(z,t) - C_0 \frac{\partial V(z,t)}{\partial t}
    \]
    \item \textbf{Lossless line:} for $R_0 = 0$ and $G_0 = 0$,
    \[
    \frac{\partial^2 V}{\partial z^2} = L_0 C_0 \frac{\partial^2 V}{\partial t^2}
    \]
\end{itemize}

\subsection*{Wave Propagation Dynamics}
\begin{itemize}
    \item \textbf{Propagation speed:}
    \[
    v = \frac{1}{\sqrt{L_0 C_0}}
    \]
    \item \textbf{Voltage solution:}
    \[
    V(z,t) = V^+(t - z/v) + V^-(t + z/v)
    \]
    \item \textbf{Characteristic impedance:}
    \[
    Z_0 = \sqrt{\frac{L_0}{C_0}}
    \]
    \item \textbf{Current waves:}
    \[
    i^+(z,t) = \frac{V^+(z,t)}{Z_0}, \qquad
    i^-(z,t) = -\frac{V^-(z,t)}{Z_0}
    \]
\end{itemize}

\subsection*{Reflections and Terminations}
\begin{itemize}
    \item \textbf{Voltage reflection coefficient:}
    \[
    \rho_v = \Gamma_L = \frac{R_L - Z_0}{R_L + Z_0}
    \]
    \item \textbf{Current reflection coefficient:}
    \[
    \rho_i = -\rho_v = \frac{Z_0 - R_L}{Z_0 + R_L}
    \]
    \item \textbf{Matched load:} if $R_L = Z_0$, then $\Gamma_L = 0$ and no reflections occur.
\end{itemize}

\subsection*{Power on the Transmission Line}
\begin{itemize}
    \item \textbf{Instantaneous power:} 
    \[
    p(z,t) = V(z,t)\, i(z,t)
    \]
    \item \textbf{Average power for a traveling wave:}
    \[
    P_{\text{avg}} = \frac{V_{\mathrm{rms}}^2}{Z_0} = I_{\mathrm{rms}}^2 Z_0
    \]
    \item \textbf{Power transfer:} maximum power is delivered when the line is matched to the load.
\end{itemize}

\section{Smatrix}
\subsection*{Scattering Parameters for an Impedance Discontinuity}
\begin{itemize}
    \item \textbf{Impedances:} Two semi-infinite lines meet at a junction with characteristic impedances $Z_1$ and $Z_2$.
    \item \textbf{Power-normalized waves:}
    \[
    a_n = \frac{V_n^+}{\sqrt{Z_n}}, \qquad b_n = \frac{V_n^-}{\sqrt{Z_n}}
    \]
    \item \textbf{Port convention:} Port 1 corresponds to $Z_1$, and Port 2 corresponds to $Z_2$.
\end{itemize}

\subsubsection*{Reflection and Transmission Coefficients}
\begin{itemize}
    \item \textbf{Reflection at the junction from Port 2:}
    \[
    S_{22} = \Gamma_{A^+} = \frac{Z_1 - Z_2}{Z_1 + Z_2}
    \]
    \item \textbf{Voltage transmission coefficient:}
    \[
    \frac{V_A^-}{V_A^+} = \frac{2Z_1}{Z_1 + Z_2}
    \]
    \item \textbf{Power-normalized transmission:}
    \[
    S_{12} = \sqrt{\frac{Z_2}{Z_1}} \cdot \frac{V_A^-}{V_A^+}
    = \frac{2\sqrt{Z_1 Z_2}}{Z_1 + Z_2}
    \]
    \item \textbf{Symmetric result:} by reciprocity,
    \[
    S_{21} = \frac{2\sqrt{Z_1 Z_2}}{Z_1 + Z_2}
    \]
    \item \textbf{Reflection seen from Port 1:}
    \[
    S_{11} = \frac{Z_2 - Z_1}{Z_1 + Z_2}
    \]
\end{itemize}

\subsubsection*{Scattering Matrix}
\[
\mathbf{S} = \frac{1}{Z_1 + Z_2}
\begin{pmatrix}
Z_2 - Z_1 & 2\sqrt{Z_1 Z_2} \\
2\sqrt{Z_1 Z_2} & Z_1 - Z_2
\end{pmatrix}
\]

\subsubsection*{Device Properties}
\begin{itemize}
    \item \textbf{Reciprocal:} $S_{12} = S_{21}$.
    \item \textbf{Not symmetric:} generally $S_{11} \neq S_{22}$ unless $Z_1 = Z_2$.
    \item \textbf{Lossless junction:} the S-matrix is unitary, $\mathbf{S}\mathbf{S}^\dagger = \mathbf{I}$.
\end{itemize}


\subsection*{Fundamental Power Wave Definitions}
\begin{itemize}
    \item \textbf{Incident wave:} 
    \[
    a_i = \frac{V_i + Z_{ri} I_i}{2\sqrt{Z_{ri}}}
    \]
    \item \textbf{Reflected wave:}
    \[
    b_i = \frac{V_i - Z_{ri} I_i}{2\sqrt{Z_{ri}}}
    \]
    \item \textbf{Power normalization:} $|a_i|^2$ and $|b_i|^2$ represent power in watts.
    \item \textbf{Net real power at port $i$:}
    \[
    P_i = |a_i|^2 - |b_i|^2
    \]
    \item \textbf{Spatial propagation:}
    \[
    a(x) = a(0)e^{-jkx}, \qquad b(x) = b(0)e^{+jkx}
    \]
    \item \textbf{Reference impedance:} $Z_{ri}$ is the characteristic impedance used to normalize port waves.
\end{itemize}

\subsection*{General N-Port S-Matrix Framework}
\begin{itemize}
    \item \textbf{Black-box model:} internal circuit details are replaced by input-output relations between incident and reflected waves.
    \item \textbf{S-matrix relation:}
    \[
    \mathbf{b} = \mathbf{S}\mathbf{a}
    \]
    \item \textbf{Matrix size:} for an $N$-port network, $\mathbf{S}$ is an $N \times N$ matrix.
    \item \textbf{Scattering parameter definition:}
    \[
    S_{ij} = \left.\frac{b_i}{a_j}\right|_{a_k = 0 \ \text{for all}\ k \neq j}
    \]
    \item \textbf{Matched-load condition:} $a_k=0$ is realized by terminating all non-driven ports with their reference impedances.
\end{itemize}

\subsection*{Two-Port Network}
\begin{itemize}
    \item \textbf{Equations:}
    \[
    b_1 = S_{11}a_1 + S_{12}a_2, \qquad
    b_2 = S_{21}a_1 + S_{22}a_2
    \]
    \item \textbf{Matrix form:}
    \[
    \begin{bmatrix} b_1 \\ b_2 \end{bmatrix}
    =
    \begin{bmatrix}
    S_{11} & S_{12} \\
    S_{21} & S_{22}
    \end{bmatrix}
    \begin{bmatrix} a_1 \\ a_2 \end{bmatrix}
    \]
\end{itemize}

\subsection*{Physical Interpretation of S-Parameters}
\begin{itemize}
    \item \textbf{$S_{ii}$:} reflection coefficient seen looking into port $i$ with all other ports matched.
    \item \textbf{$S_{ij}$:} transmission coefficient from port $j$ to port $i$.
    \item \textbf{$S_{21}$:} forward transmission (gain/loss).
    \item \textbf{$S_{12}$:} reverse isolation.
    \item \textbf{Power directionality:} S-parameters quantify how much power is reflected or transmitted through the network.
\end{itemize}

\section{VNA}

\subsection*{Power Wave Variable Definitions}
\begin{itemize}
    \item \textbf{$a_1$:} incident wave at Port 1.
    \item \textbf{$b_1$:} power reflected from Port 1.
    \item \textbf{$b_2$:} power transmitted through the DUT and leaving Port 2.
    \item \textbf{$a_2$:} incident wave at Port 2.
\end{itemize}

\subsection*{Directional Coupler Performance}
\begin{itemize}
    \item \textbf{Coupling factor (dB):}
    \[
    \text{Coupling (dB)} = -10\log\!\left(\frac{P_{\text{coupling forward}}}{P_{\text{incident}}}\right)
    \]
    \item \textbf{Isolation factor (dB):}
    \[
    \text{Isolation (dB)} = -10\log\!\left(\frac{P_{\text{coupled reverse}}}{P_{\text{incident}}}\right)
    \]
    \item \textbf{Directivity (dB):}
    \[
    \text{Directivity (dB)} = \text{Isolation (dB)} - \text{Coupling (dB)}
    \]
    \item \textbf{Mainline loss:} total insertion loss in the primary through-line.
\end{itemize}

\subsection*{VNA Measurement Signal Relationships}
\begin{itemize}
    \item \textbf{Working principle:} a 3-way resistive splitter is used in simplified models to sample the input signal.
    \item \textbf{Instrument architecture:} actual VNA hardware uses directional couplers to separate incident and reflected waves.
    \item \textbf{Coherent receivers:} RX REF1, RX TEST1, RX REF2, and RX TEST2 share a common reference oscillator.
\end{itemize}

\subsection*{S-Parameter Extractions}
\begin{itemize}
    \item \textbf{Forward measurements:} when SW1 is in position 1,
    \[
    S_{11} = \frac{b_1}{a_1}, \qquad S_{21} = \frac{b_2}{a_1}
    \]
    \item \textbf{Reverse measurements:} when SW1 is in position 2,
    \[
    S_{22} = \frac{b_2}{a_2}, \qquad S_{12} = \frac{b_1}{a_2}
    \]
    \item \textbf{2$\times$2 S-matrix:}
    \[
    \begin{bmatrix}
    S_{11} & S_{12} \\
    S_{21} & S_{22}
    \end{bmatrix}
    \]
    \item \textbf{Index convention:} $i$ = destination/output port, $j$ = source/input port.
\end{itemize}

\subsection*{Laboratory Reference}
\begin{itemize}
    \item \textbf{Core function:} characterizes amplitude and phase response over frequency without knowing the DUT topology.
    \item \textbf{Applications:} filters, amplifiers, antennas, resonators, and couplers.
    \item \textbf{Systems engineering:} with antennas, the VNA can function as a radar system.
    \item \textbf{Advanced diagnostics:} can support material characterization and imaging applications.
\end{itemize}

\section{Modulations}

\subsection*{Reference Nomenclature and Variable Definitions}
\begin{itemize}
    \item \textbf{$f_s$:} DAC/ADC sampling rate.
    \item \textbf{$f_{rfmax}$:} upper bound of the allocated RF band.
    \item \textbf{$f_c$:} RF center frequency (carrier frequency).
    \item \textbf{$f_{C-IF}$:} local oscillator frequency for IF translation.
    \item \textbf{$f_{IF}$:} intermediate frequency.
    \item \textbf{$f_{s(BB)}$:} baseband sampling rate.
    \item \textbf{ADC:} analogue-to-digital converter.
    \item \textbf{DAC:} digital-to-analogue converter.
    \item \textbf{DSP:} digital signal processing.
\end{itemize}

\subsection*{Nyquist Sampling Criterion for RF Systems}
\begin{itemize}
    \item \textbf{Primary condition:}
    \[
    f_s > 2f_{rfmax}
    \]
    \item \textbf{Condition met:} the digital/analogue interface can be placed directly at the RF stage.
    \item \textbf{Condition not met:} a staged architecture using an intermediate frequency $f_{IF}$ is required.
\end{itemize}

\subsection*{Architectural Definitions}
\begin{itemize}
    \item \textbf{Heterodyning:} direct translation between baseband and RF in a single step.
    \item \textbf{Superheterodyning:} multi-stage translation using an intermediate frequency $f_{IF}$.
    \item \textbf{Stage 1:} baseband to intermediate frequency.
    \item \textbf{Stage 2:} intermediate frequency to radio frequency.
\end{itemize}

\subsection*{Comparative Analysis of Radio Architectures}
\begin{itemize}
    \item \textbf{Direct RF:} ADC/DAC at the RF stage near the antenna; requires very high sampling rates.
    \item \textbf{IF Sampling:} digital processing occurs at $f_{IF}$, while analogue mixers and LOs perform the RF translation.
    \item \textbf{Baseband Sampling:} the interface is placed at baseband, and modulation/demodulation occurs in analogue hardware.
    \item \textbf{Interface position logic:} higher sampling rates move the digital/analogue interface closer to the antenna.
\end{itemize}

\subsection*{Technical Advantages of Digital Implementation}
\begin{itemize}
    \item \textbf{Operational precision:} digital logic provides deterministic and accurate operation.
    \item \textbf{Hardware robustness:} reduced drift from tolerances, thermal effects, and ageing.
    \item \textbf{System efficiency:} smaller footprint, simpler BOM, and lower total power consumption.
    \item \textbf{Flexibility:} reconfigurable hardware supports multiple communication standards via software.
\end{itemize}

\subsection*{Summary Reference Table}
\begin{center}
\begin{tabular}{lll}
\textbf{Architecture} & \textbf{Typical Sampling Rate} & \textbf{Interface Position} \\
\hline
Direct RF & $\approx 10$ GSps & RF stage near antenna \\
IF Sampling & 10s MSps & Intermediate frequency $f_{IF}$ \\
Baseband Sampling & 10s--100s kSps & Baseband stage
\end{tabular}
\end{center}

\section{Heterodyning}

\subsection*{Fundamental Sampling Criterion}
\begin{itemize}
    \item \textbf{Nyquist-Shannon condition:}
    \[
    f_s > 2f_{rfmax}
    \]
    \item \textbf{$f_s$:} sampling rate of the ADC or DAC.
    \item \textbf{$f_{rfmax}$:} maximum frequency component in the RF-modulated signal.
    \item \textbf{Consequence:} if the condition is not met, an intermediate frequency stage is required.
\end{itemize}

\subsection*{Modulation and Demodulation Schemes}
\begin{itemize}
    \item \textbf{Heterodyning:} single-stage translation between baseband and RF using carrier frequency $f_c$.
    \item \textbf{Superheterodyning:} two-stage translation via intermediate frequency $f_{IF}$.
    \item \textbf{Stage 1:} baseband $\rightarrow$ IF.
    \item \textbf{Stage 2:} IF $\rightarrow$ RF.
    \item \textbf{Typical IF band:} 10s of MHz to 100s of MHz.
    \item \textbf{Typical RF band:} 100s of MHz to 10s of GHz.
\end{itemize}

\subsection*{Comparative Analysis of Radio Architectures}
\begin{center}
\begin{tabular}{llll}
\textbf{Architecture} & \textbf{Interface Placement} & \textbf{Typical Sampling Rate} & \textbf{Logic} \\
\hline
Direct RF & RF stage & $\approx 10$ GSps & DUC/DDC with interpolators, decimators, and NCOs \\
IF Sampling & IF stage & 10s MSps & Analog IF-to-RF translation with image rejection filters \\
Baseband Sampling & Baseband stage & 10s--100s kSps & Analog-domain modulation/demodulation
\end{tabular}
\end{center}

\subsection*{Digital Processing Advantages}
\begin{itemize}
    \item \textbf{Operational precision:} deterministic accuracy and reduced nonlinear errors.
    \item \textbf{Long-term reliability:} reduced sensitivity to drift and aging effects.
    \item \textbf{Physical and fiscal efficiency:} smaller footprint, simpler BOM, and lower power.
    \item \textbf{Architectural versatility:} reconfigurable hardware for multiple waveforms and bands.
\end{itemize}

\subsection*{RFSoC Integration}
\begin{itemize}
    \item \textbf{RFSoC:} integrates high-performance ADCs/DACs, digital logic, and processor cores on one CMOS die.
    \item \textbf{Benefit:} enables direct digital modulation with reduced analog complexity.
\end{itemize}

\end{document}
